\documentclass[11pt]{extarticle}
\usepackage{manualdoprofessor}
\usepackage{fichatecnica}
\usepackage{lipsum,media9}
\usepackage[justification=raggedright]{caption}
\usepackage[one]{bncc}
\usepackage[acorde]{../edlab}
\usepackage{marginnote}
\usepackage{pdfpages}
\usepackage[printwatermark]{xwatermark}
% \newwatermark[pagex=2]{\includegraphics[scale=3.3]{watermarks/test-a.png}}	% página específica
% %\newwatermark[oddpages]{\includegraphics{watermarks/test-a.png}}			% páginas ímpars
% %\newwatermark[evenpages]{\includegraphics{watermarks/test-a.png}}			% págimas pares
% \newwatermark[allpages]{\includegraphics[scale=3.3]{watermarks/test-b.png}}
% \pagecolor{cyan!0!magenta!10!yellow!28!black!28!}

\newcommand{\AutorLivro}{Arthur Nestrovski}
\newcommand{\TituloLivro}{Os doze trabalhos de Hércules}
\newcommand{\Genero}{Lendas, mitos, fábula}
%\newcommand{\imagemCapa}{./images/PNLD0001-01.png}
\newcommand{\issnppub}{978-65-99441-24-0}
\newcommand{\issnepub}{978-65-99441-27-1}
% \newcommand{\fichacatalografica}{PNLD0001-00.png}
\newcommand{\colaborador}{Gabriela Karam}

\begin{document}

\title{\TituloLivro}
\author{\AutorLivro}
\def\authornotes{\colaborador}

\date{}
\maketitle

%\begin{abstract}\addcontentsline{toc}{section}{Carta ao professor}
%\pagebreak

\tableofcontents


\begin{abstract}
Este material tem a intenção de contribuir para que você desenvolva um trabalho aprofundado com a obra \textit{Os doze trabalhos de Hércules} em sala de aula. Você encontrará informações sobre o autor, sobre o gênero e também algumas propostas de trabalho para a sala de aula que você poderá explorar livremente, da forma que considerar mais apropriada para os seus estudantes. 

A principal razão que leva uma pessoa a estudar os mitos é a compreensão do ser humano. O fato é que, no mundo racional da sociedade contemporânea, nos afastamos cada vez mais dos sonhos e não entendemos bem o papel dos mitos, dos heróis e das cerimônias rituais, estas, reguladoras das práticas sociais. Perdemos o sentido das metáforas. Não sabemos desvelá-las e o ponto de vista cartesiano nos faz acreditar que essas histórias são primitivas e não nos dizem respeito. O autor Joseph Campbell -- uma das maiores autoridades mundiais nesse assunto, em anos de pesquisas sobre os mitos de inúmeras sociedades -- formulou ideias originais sobre a similaridade entre os povos em suas relações com o cosmos. Tais noções nos convidam a buscar uma nova forma de interpretar a nossa gênese. Em nossa diversidade, ele ressalta que somos um único povo – o povo terrestre – pois do alto, do espaço sideral, a Terra é vista sem fronteiras políticas, sem barreiras que impeçam a nossa inter-relação. 

Em \textit{Os doze trabalhos de Hércules}, o autor Arthur Nestrovski conta, de forma leve e despretensiosa, um mito grego muito conhecido, com o apoio de ilustrações de Zansky, que complementam essa aventura de Héracles, um semideus filho de Zeus e de Alcmena, filha de um rei mortal. Portanto, Herácles é um dos filhos que Zeus teve fora do casamento com Hera. Héracles se tornou o mais forte e poderoso de todos os mortais, uma espécie de super-herói mitológico. Os romanos o chamavam de Hércules, nome como ficou mais conhecido com o passar dos tempos. No mito original, o semideus Hércules teve de cumprir doze tarefas como penitência por ter, sob o efeito de um feitiço, matado a mulher e os filhos. 

Na obra, dirigida ao público infantil, Hera, mulher de Zeus, imbuída de ciúmes, consegue transformar seu menino humano predileto, Euristeu, em rei, e obriga Hércules a se submeter a ele caso queira garantir a imortalidade. Foi por conta disso que Hércules teve de realizar os famosos doze trabalhos, ordenados pelo rei Euristeu. 

Ao longo do manual, todos esses aspectos serão explorados e relacionados a sugestões de atividades. Com isso, nossa proposta é oferecer ideias e inspirações para um trabalho que pode ser desenvolvido tanto a curto, quanto a médio e longo prazo. Sinta-se à vontade para personalizar a aula e torná-la sua, aplicando seus conhecimentos, sua personalidade e aproveite para fortalecer seu vínculo com a turma. Boa aula! 

\end{abstract}

\section{Sobre o livro}
\Image{Os deuses gregos (Pinterest; Domínio público)}{PNLD2023-034-02.png}

\section{Sobre o autor}

\paragraph{O autor} Arthur Nestrovski nasceu em Porto Alegre, no ano de 1959. Desde 2010 é diretor artístico da Orquestra Sinfônica do Estado de São Paulo, a OSESP, maior orquestra do Brasil. É Doutor em literatura e música pela Universidade de Iowa (EUA) e segue carreira musical, no Brasil e no exterior, como violonista e compositor. É parceiro de grandes compositores brasileiros, como José Miguel Wisnik, Zélia Duncan e Adriana Calcanhotto. 

Além de \textit{Os Doze Trabalhos de Hércules}, Nestrovski escreveu outros dez livros para crianças, entre eles \textit{Histórias de avô e avó} (Cia. das Letrinhas, 1998), \textit{Bichos que existem \& bichos que não existem} (Cosac
Naify, 2002, vencedor do Prêmio Jabuti de ``Livro do Ano'')
e \textit{Pelo nariz} (Cosac Naify, 2013,“Melhor Livro Infantil'' da Biblioteca Nacional). 

\paragraph{O ilustrador} Zansky nasceu em 1979 e é ilustrador desde 2008. É bacharel em Artes 
Plásticas pela Unesp. Trabalha com livros, revistas, vídeos, 
capas e encartes de álbuns e posters.  


\section{Sobre o gênero}

\paragraph{O gênero} O gênero deste livro é \textit{Lendas; mitos; fábula}. 

A lenda e o mito são narrativas fantasiosas transmitidas pela tradição oral através dos tempos. De caráter fantástico, as lendas e os mitos combinam fatos reais e históricos com eventos que não têm comprovação de acontecimento, a não ser pela palavra dos que sobraram para contar a história. As lendas e mitos de uma sociedade são fundamentais para que entendamos quem são essas pessoas e no que acreditam, bem como suas tradições. Uma lenda é verdadeira até que se prove o contrário. Com exemplos bem definidos em todos os países do mundo, as lendas e os mitos de um povo geralmente fornecem explicações plausíveis, e até certo ponto aceitáveis, para coisas que não têm explicações científicas comprovadas, como acontecimentos misteriosos ou sobrenaturais.

A fábula é uma narrativa curta em que os personagens principais geralmente são seres personificados. Esses seres apresentam características humanas, tais como a fala e traços de personalidade. Essas personagens podem ser também objetos animados ou deuses. Em cada história há uma lição de moral: uma mensagem de cunho educativo que busca conscientizar o leitor. A fábula tem estreita relação com o gênero conto, mas se diferencia pela centralidade dos personagens animais e pelo intuito de concluir a história com um ensinamento. É uma história que pode ser contada em prosa ou em versos. 

Sobre a origem da fábula, Douglas Tufano afirma que:

\begin{quote}
A fábula teria nascido provavelmente na Ásia Menor e daí teria passado pelas ilhas gregas, chegando ao continente helênico. Há registros sobre fábulas egípcias e hindus, mas sua criação é atribuída à Grécia, pois é onde a fábula passa a ser considerada como um tipo específico de criatividade dentro da teoria literária. 

Na Grécia, os primeiros exemplos de fábula datam do século \textsc{viii} a.C. Isso nos mostra, é claro, que Esopo não foi o inventor do gênero, mas sim o mais conhecido fabulista na Antiguidade como autor e narrador dessas pequenas histórias.\footnote{\textsc{tufano}, Douglas. \textit{Esopo: Fábulas completas}. São Paulo: Moderna, 2015.}
\end{quote}

Esopo foi um autor da Grécia Antiga a quem são atribuídas algumas das mais famosas fábulas, como \textit{A raposa e o cacho de uvas} e \textit{A galinha de ovos de ouro}. Diversas  histórias suas foram recontadas por La Fontaine, que é também um dos mais clássicos fabulistas do Ocidente.

No caso de \textit{O reino do meio da tarde}, há os elementos típicos das lendas e fábulas.
A história se inicia com o clássico ``era uma vez'', que já situa a narrativa no universo fantástico das lendas, estrutura típica das fábulas e histórias da carochinha. A descrição que se segue é de um reino mágico, caracterizado por uma condição especial (um lugar onde não havia tempo), seus habitantes fantásticos, que são os elementos do mundo personificados (ferros de passar alados, funis triunfantes, a princesa Lua etc.) e o comportamento particular de seus astros e planetas (sóis que habitam ninhos e estrelas ninadas em berçários de constelações).

A estrutura do enredo também segue o padrão das narrativas orais que caracteriza as lendas e fábulas: há uma situação inicial dada desde o início dos tempos (o reino bonito e tranquilo, escondido no meio das nuvens); o surgimento de um problema a ser contornado (os súditos que se cansaram de viver apenas na tarde e queriam o tempo em seu reino); as tentativas para solucionar o problema (as expedições dos cidadãos do reino em busca da manhã e da noite); o aparecimento do herói que contorna o problema (a princesa Lua que consegue trazer uma estrela ao reino e, assim, inaugurar a manhã e a noite); e o desfecho feliz, no qual o herói retorna a casa e o problema é solucionado (todos no reino ficam felizes, pois passaram a ter o tempo que tanto almejaram).

\section{Proposta de Atividades}
\subsection{Pré Leitura}

\BNCC{EF15LP02}% Estabelecer expectativas em relação ao texto que vai ler (pressuposições antecipadoras dos sentidos, da forma e da função social do texto), apoiando-se em seus conhecimentos prévios sobre as condições de produção e recepção desse texto, o gênero, o suporte e o universo temático, bem como sobre saliências textuais, recursos gráficos, imagens, dados da própria obra (índice, prefácio etc.), confirmando antecipações e inferências realizadas antes e durante a leitura de textos, checando a adequação das hipóteses realizadas.}
\BNCC{EF03LP26}% Identificar e reproduzir, em relatórios de observação e pesquisa, a formatação e diagramação específica desses gêneros (passos ou listas de itens, tabelas, ilustrações, gráficos, resumo dos resultados), inclusive em suas versões orais.
\BNCC{EF04HI02}% Identificar mudanças e permanências ao longo do tempo, discutindo os sentidos dos grandes marcos da história da humanidade (nomadismo, desenvolvimento da agricultura e do pastoreio, criação da indústria etc.)
\BNCC {EF04ER04}% Identificar as diversas formas de expressão da espiritualidade (orações, cultos, gestos, cantos, dança, meditação) nas diferentes tradições religiosas.

\paragraph{Tema} Reflexões iniciais sobre mitologia. 

\paragraph{Conteúdo} Mitologia Grega.

\paragraph{Objetivo} Introduzir os estudantes no universo da Mitologia Grega. 

\paragraph{Justificativa} O mundo da tecnologia e da produtividade afasta do cotidiano, cada vez mais, o universo das lendas e dos mitos. Estudar esse universo é, primeiramente, uma forma de imergir numa outra forma de conhecimento e estimular a imaginação dos estudantes. Conhecer a mitologia é entender-lhe a função social e a religiosidade, analisando-as à luz do mundo contemporâneo.  

\paragraph{Metodologia}

Antes de mergulhar no universo de \textit{Os Doze Trabalhos de Hércules}, o educador pode estimular a turma a explorar, em outras fontes de informação, a mitologia de uma forma geral. Para isso, incentive a pesquisa, em livros e sites de internet, sobre mitologia. A indicação inicial para promover uma problematização do gênero da obra em questão é o vídeo \href{https://youtu.be/9qlhKoUj_dw}{\textit{A mitologia grega é adequada para crianças?}}, do canal \textit{Noites Gregas}, no qual o professor Moreno problematiza a mitologia grega para crianças. É um ponto de vista que nos faz pensar sobre o assunto e pode gerar bons debates com as crianças. O professor Moreno desenvolve um pensamento a respeito das modificações que são feitas nos mitos para que o público infantojuvenil possa também ter um primeiro contato com essas histórias. 

Proponha que as crianças se sentem em forma de círculo com suas carteiras para que aconteça um debate acerca do que esperam das histórias mitológicas, auxiliando-as a pensar nos pontos interessantes que esse tipo de obra pode trazer para a vida deles, e quais as visões que elas têm a respeito dos mitos. Depois do debate, o professor pode sugerir que todos façam uma pesquisa em conjunto sobre alguns pontos:

\begin{itemize}
\item O que é a mitologia grega e o que ela significa para o ser humano mesmo depois de tanto tempo?
\item Quais as ligações entre os mitos e o mundo contemporâneo?
\item Quais eram as características históricas do mundo grego?
\item Onde está a Grécia no mapa e que características a levaram para a criação de mitos?
\item Qual era a relação entre a mitologia grega e a religiosidade desse povo?
\item Quais as diferenças entre o politeísmo, encontrado na Grécia Antiga, e o monoteísmo? %PNLD2023-034-04
\end{itemize}

\Image{A mitologia grega (Pinterest; Domínio público)}{PNLD2023-034-03.png}

\SideImage{O minotauro (Pinterest; Domínio público)}{PNLD2023-034-04.png}

Para auxiliar nessas reflexões, sugerimos iniciar pelo \url{https://youtu.be/20bVidCkrdw}{vídeo que conta o mito do Minotauro} e depois o da \url{https://youtu.be/6EG5fEECKYY}{Medusa}. Logo depois, para aproximar ainda mais os alunos do universo da mitologia, é interessante fazer com que a turma conheça mais os deuses e deusas gregos. Sugerimos a leitura e o debate sobre essas personagens, para que as crianças consigam entender a importância e a grandeza dessas personagens imortais e poderosas. Logo em seguida, pode ser mostrada a animação cantada \textit{Os deuses do olimpo}, do canal \href{https://youtu.be/DJ-5cxJVI64}{\textit{Um pouco sobre quase tudo}}, que conta a história dos deuses do olimpo, entre eles Zeus -- pai da nossa personagem Hércules. 

\Image{O professor pode trazer uma série de repertório (Educa Mais Brasil; Domínio público)}{PNLD2023-034-05.png}

Ao a exibição dos vídeos e da pesquisa em si, realize junto com a turma um relatório sobre as informações que colherem, baseado nos itens norteadores, levando em conta as características principais dos mitos gregos, suas funções sociais, a presença da religiosidade e a conexão de tudo isso com o mundo contemporâneo. Essas e outras questões podem ser debatidas em sala de aula, assim como o significado da metáfora e como ela é usada no contexto mitológico, sempre dando liberdade para que os alunos expressem dúvidas e questionamentos.

\SideImage{Realize junto com a turma um relatório sobre as informações que colherem (Educa Mais Brasil; Domínio público)}{PNLD2023-034-06.png}

\paragraph{Tempo Estimado} Quatro aulas de 50 minutos.

\subsection{Leitura}

\BNCC{EF35LP21}% Ler e compreender, de forma autônoma, textos literários de diferentes gêneros e extensões, inclusive aqueles sem ilustrações, estabelecendo preferências por gêneros, temas, autores.
\BNCC{EF35LP03}% Identificar a ideia central do texto, demonstrando compreensão global.
\BNCC{EF15LP16}% Ler e compreender, em colaboração com os colegas e com a ajuda do professor e, mais tarde, de maneira autônoma, textos narrativos de maior porte como contos (populares, de fadas, acumulativos, de assombração etc.) e crônicas.

\paragraph{Tema} Leitura e compreensão de \textit{Os Doze Trabalhos de Hércules}.

\paragraph{Conteúdo} Estrutura de de \textit{Os Doze Trabalhos de Hércules}.  

\paragraph{Objetivo} Ler e compreender o texto de \textit{Os Doze Trabalhos de Hércules} e seus pressupostos, especialmente os conhecimentos a respeito da mitologia, estudados na atividade de Pré-Leitura. Promover o gosto pelos livros e pela leitura, estimular o desenvolvimento da linguagem e da criatividade.   

\paragraph{Justificativa} O professor tem grande influência na formação de leitores, especialmente quando se trata de mitos e lendas, para ampliar o repertório cultural dos alunos. É preciso fazer a mediação entre a obra, sua linguagem, suas estruturas, seus pressupostos e os estudantes, de preferência estabelecendo uma relação fundamentada no prazer, na identificação e na liberdade de interpretação. Eis o nosso desafio: ler com os alunos, apresentando as passagens decisivas de um texto, explicando por que elas chamam a atenção e  ouvindo as impressões dos estudantes a respeito de tudo isso. 

Após a pesquisa sobre mitologia e um primeiro contato com Hércules, nosso herói semideus, divida a turma em grupos responsáveis por cada um dos trabalhos de Hércules. Cada coletivo deve ler a sua parte -- escolhida pelo professor -- e preparar uma narrativa para contá-la para toda a turma. Sugerimos que o educador inicie os trabalhos lendo as primeiras páginas do livro, que explica o nascimento do herói e como ele foi submetido aos doze trabalhos. Logo em seguida, organize a classe de modo que cada grupo se prepare para expor a sua narração. Quando a leitura acabar, em círculo, estimule os estudantes a debaterem sobre qual narração acharam mais interessante. A partir daí proponha um debate sobre a oralidade das histórias e a importância de transmiti-las de modo que todas as pessoas entendam e possam repassá-las à sua maneira. 

Ainda em roda, proponha que os estudantes contem histórias que eles já tenham ouvido dos seus pais, avós, tias e tios, e que possam ser classificadas como mitos. Relacione essas narrativas com elementos que foram vistos na obra, especialmente as metáforas. 

\Image{Realize junto com a turma um relatório sobre as informações que colherem (Educa Mais Brasil; Domínio público)}{PNLD2023-034-07.png}

\paragraph{Tempo Estimado} Quatro aulas de 50 minutos.

\subsection{Pós-leitura}

\BNCC{EF35LP25}% Criar narrativas ficcionais, com certa autonomia, utilizando detalhes descritivos, sequências de eventos e imagens apropriadas para sustentar o sentido do texto, e marcadores de tempo, espaço e de fala de personagens.
\BNCC{EF05HI08}% Identificar formas de marcação da passagem do tempo em distintas sociedades, incluindo os povos indígenas originários e os povos africanos.}
\BNCC{EF05HI09}% Comparar pontos de vista sobre temas que impactam a vida cotidiana no tempo presente, por meio do acesso a diferentes fontes, incluindo orais.}

\paragraph{Tema} Mitologia. 

\paragraph{Conteúdo} Comparação entre mitologia grega e indígena.

\paragraph{Objetivo} Comparar diferentes culturas. Desenvolver exercícios de escrita.

\paragraph{Justificativa} Enquanto as atividades de leitura, compreensão e análise caracterizam-se pelas primeiras aproximações do texto, seguidas de atividades de descrição de suas características, a prática de redação abre espaço para que os estudantes criem suas próprias narrativas. Além disso, a comparação entre mitos gregos e indígenas serve para a compreensão da diferença entre essas culturas. 

\paragraph{Metodologia} Após a leitura, é interessante fazer um comparativo entre as expectativas que foram relatadas no primeiro momento da atividade com o que foi lido e aprendido na obra. Proponha um debate em que os estudantes articulem suas opiniões. Depois de uma breve discussão, inicie uma pesquisa sobre outras mitologias além das existentes no mundo grego. Compare as histórias narradas na atividade anterior com outras, de nossas culturas originárias, sejam elas africanas ou indígenas.  

Para essas atividades, a matéria \href{https://super.abril.com.br/mundo-estranho/quais-sao-os-principais-deuses-da-mitologia-indigena-brasileira/}{\textit{Quais são os principais deuses da mitologia indígena brasileira}}, da \textit{Revista super interessante}, pode nortear e enriquecer a discussão. O artigo traz alguns dos deuses da mitologia indígena, comentando suas características e descrevendo os mitos que fundam essas figuras, bem como a noção da passagem do tempo para esses povos. 

Auxilie os estudantes a debaterem sobre as diferenças entre a mitologia grega e a indígena, buscando relacionar esses mitos com o tempo presente. Sugira também que os alunos, divididos em novos grupos, redijam em conjunto um mito com base no livro e em tudo o que foi pesquisado até agora. 

Promova uma busca sobre o que mais poderia ser contado para que essa história cause interesse para quem ouve, para que ela possa passar de geração a geração e até ser transformada em livro. Estimule a criação de um livro dos mitos que foram criados pelos grupos, levando em consideração os comentários feitos sobre a narrativa, pedindo, ainda, que haja uma imagem que represente esse mito. Ela pode ser feita usando diferentes materiais, como lápis de cor, giz de cera, tinta guache e o que mais a escola disponibilizar. 

Sugerimos que todos os mitos, criados por cada grupo, deem origem a  livros que ficarão na biblioteca da escola para que possam ser lidos e passados adiante assim como ocorre em grande parte das culturas, tanto ocidentais como orientais.

\paragraph{Tempo Estimado} Quatro aulas de 50 minutos.

\Image{Após a imersão, sugira que os alunos, divididos em novos grupos, usem muita criatividade e redijam em conjunto um mito com base em tudo o que foi pesquisado e na leitura (Mitologia; Domínio público)}{PNLD2023-034-08.png}


\section{Sugestões de referências complementares}

\subsection{Livros} 

\begin{itemize}
\item \textsc{campbell}, Joseph. \textit{O poder do mito}. São Paulo: Cultrix, 1990.
\SideImage{Realize junto com a turma um relatório sobre as informações que colherem (Abril; Domínio público)}{PNLD2023-034-09.png}

Livro proveniente de uma série de conversas mantidas entre Joseph Campbell e o jornalista Bill Moyers a respeito de mitologia.
%PNLD2023-034-10
\end{itemize}

\section{Bibliografia comentada}
\subsection{Livros}

\begin{itemize}
\item \textsc{brasil}. Ministério da Educação. \textit{Base Nacional Comum Curricular}. Brasília, 2018.

Consultar a \textsc{bncc} é essencial para criar atividades para a turma. Além de especificar quais habilidades precisam ser desenvolvidas em cada ano, é fonte de informações sobre o processo de aprendizagem infantil. 
 
 \item \textsc{grimm}, Jacob e Wilhelm. \textit{Contos maravilhosos infantis e domésticos}. São Paulo: Editora 34, 2018.

Inúmeros contos fantásticos que foram publicados inicialmente em 1812. Contos clássicos que originaram diferentes histórias conhecidas no mundo ocidental.

\item \textsc{kimmel}, Eric. \textit{Mitos gregos}. São Paulo: WMF Martins Fontes, 2013.

Livro em que o autor se preocupa em aproximar os heróis gregos do jovem leitor de hoje.

\item \textsc{munduruku}, Daniel. \textit{Contos indígenas brasileiros}. São Paulo: Global Editora, 2004.

Livro que traz contos dos povos originários e apresenta a diversidade cultural e linguística no Brasil.

\item \textsc{coelho}, Nelly Novaes. \textit{Literatura infantil, teoria, análise, didática}. 1ª ed. São Paulo: Moderna, 2000.

Livro que fala sobre os espaços da literatura infantil na contemporaneidade e a importância de as crianças estarem ligadas ao seu imaginário pela via literária.

\end{itemize}

\subsection{\textit{Sites}}

\begin{itemize}
\item Artigo \href{https://siteantigo.portaleducacao.com.br/conteudo/artigos/educacao/a-importancia-da-leitura-dos-contos-de-fadas-na-educacao-infantil/30151}{``A Importância da Leitura dos Contos de Fadas na Educação Infantil''}, por Ana Maria da Silva.  


No artigo, a autora comenta a importância da construção do imaginário por meio da literatura para as crianças, trazendo elementos que analisam o mundo pós-moderno e os espaços que a literatura infantil, principalmente os contos, devem ter.

\item Artigo \href{http://docs.uninove.br/arte/fac/publicacoes/pdf/v3-n1-2012/Francy.pdf}{``Literatura infantil: A contribuição dos contos de fadas para a construção do imaginário infantil''}.

\item Artigo \href{https://danielmcarlos.wordpress.com/2014/02/09/mitologia-para-criancas/}{``Mitologia para crianças''}, do blog ``Filosofia animada''. 

Artigo que se propõe a mostrar, de forma lúdica, parte da mitologia grega para o público infantil.

\end{itemize}

\subsection{\textit{Filmes}}

\begin{itemize}
\item \textit{Hércules}. Dirigido por Ron Clements e John Musker, 1997.

Meio homem, meio deus, o semideus Hércules anda atormentado pela terra. Depois de doze árduos trabalhos realizados e da perda de sua família, ele conhece seis assassinos impiedosos e une-se ao grupo em busca de trabalho, até que o rei da Trácia convida-o para treinar o seu exército.

\item \textit{Fúria de titãs}. Dirigido por Louis Leterrier, 2010.

Perseu descobre que é o filho mortal de Zeus, mas recusa-se a aceitar tal condição. Ele encontra-se entre a batalha dos deuses e sem ajuda para salvar a cidade e sua família da vingança de seu tio Hades, o deus do submundo. Sem nada a perder, Perseu lidera um grupo de guerreiros em uma perigosa missão para prevenir que Hades cause uma devastação na Terra.

\item \textit{Percy Jackson e o Ladrão de Raios}. Dirigido por Chris Columbus, 2010.

A vida do adolescente Percy Jackson, que está sempre pronto para entrar em uma confusão, torna-se bem mais complicada quando ele descobre que é filho do deus grego Poseidon. Em um campo de treinamento para filhos das divindades, Percy aprende a tirar proveito de seus poderes divinos e se preparea para a maior aventura de sua vida.

\end{itemize}
\end{document} 
